%!TEX root = ../thesis.tex
\chapter{\implementation}
\label{ch:implementation}

\begin{lstlisting}[language=erlang, frame=single, breaklines=true, numbers=left, caption={Hello World Program in Erlang.}]
-module(hello).
-export([hello_world/0]).

hello_world() -> io:fwrite("hello, world\n").
\end{lstlisting}

\begin{lstlisting}[language=cpp, frame=single, breaklines=true, numbers=left, caption={Program in C++.}]
#include <iostream>
int main() {
    std::cout << "Hello, World!" << std::endl;
    this->run();
    int a = 0;
    if (a <= 5){
        return 1;
    }
    // comment line
    return 0;
}
\end{lstlisting}

\begin{lstlisting}[language=yaml, frame=single, breaklines=true, numbers=left, caption={YAML Example}]
# A sample yaml file
key: value
array:
  - item1
  - item2:
        anotherKey: anotherValue
\end{lstlisting}

\newpage

\begin{lstlisting}[language=pseudo, frame=single, breaklines=true, numbers=left, caption={Pseudo Code Example}]
items = [1,2,3]
# Print each item in the list
for i in items:
    if i > 2:
        print("MyItem: "+i)
\end{lstlisting}



\begin{lstlisting}[language=docker, frame=single, breaklines=true, numbers=left, caption={Dockerfile Example}]
FROM python:3.13
WORKDIR /usr/local/app

# Install the application dependencies
COPY requirements.txt ./
RUN pip install --no-cache-dir -r requirements.txt

# Copy in the source code
COPY src ./src
EXPOSE 8080

# Setup an app user so the container doesn't run as the root user
RUN useradd app
USER app

CMD ["uvicorn", "app.main:app", "--host", "0.0.0.0", "--port", "8080"]
\end{lstlisting}

\newpage

\begin{lstlisting}[language=json, frame=single, breaklines=true, numbers=left, caption={JSON Example}]
{
    "key": "value",
    "pair": [
        "first": 0,
        "second": "null"
        # comment line
    ]
}
\end{lstlisting}
